\begin{beginnerstatement}[In simple terms: From version to package]

A version is a unique designation for a particular state; in the template,
\texttt{VERSION} is the central source for it. A release is a version approved
for publication, but the workflow examined here does not itself publish
anything. \texttt{release check} merely checks the prerequisites and produces
no artifacts. During \emph{packaging}, program files are combined into a
deliverable package or \emph{bundle}; this result is an artifact. A Git tag
marks a particular commit with a name, often the version number. During
\emph{signing}, a digital signature backed by a certificate confirms a
package's origin; notarization is an additional external check, and an updater
distributes later versions to installed applications. Deployment means making
an application available; Docker builds images for this purpose, containers
run them, and Docker Compose describes multiple related containers. AWS,
Azure, and GCP are possible cloud providers, but this provider-agnostic
foundation does not automatically commit to any of them.

\end{beginnerstatement}
